\section{Experiments}
\label{sec:experiments}
\label{sec:dataset}
We show that geometric calibration is biased by its feature detections
(\cref{sec:detector-bias}) and that the radiometric camera improves
the geometric accuracy (\cref{sec:geometric-accuracy}), compare the
learned vignetting with a measured flat field
(\cref{sec:experiment-vignetting}) and validate the focus model
(\cref{sec:experiment-defocus}). All captures are RAW images of a Sony
ILME-FX3 ($4256\times2848$ pixels) and are listed in
\cref{sec:supp-dataset} together with the runtime for each experiment; all errors are in pixels of this image. 

% Table of Sec. 5.2 (Geometric Accuracy), declared here so that it floats to the top of the page where the experiments start.
\begin{table*}[t]
        \centering
        \footnotesize
\setlength{\tabcolsep}{3pt}
\begin{tabular*}{\linewidth}{@{\extracolsep{\fill}}lcccccc@{}}
\toprule
 & \multicolumn{2}{c}{\textbf{Tamron 20\,mm}} & \multicolumn{2}{c}{\textbf{Sigma 14\,mm}} & \multicolumn{2}{c}{\textbf{Sigma 8\,mm fisheye}} \\
\cmidrule(lr){2-3}\cmidrule(lr){4-5}\cmidrule(lr){6-7}
\textbf{Model} & AprilGrid & Siemens stars & AprilGrid & Siemens stars & AprilGrid & Siemens stars \\
\midrule
\multicolumn{7}{@{}l}{\underline{\textit{Median reprojection error [px]}}} \\
Brown--Conrady & 0.119\,$\rightarrow$\,\textbf{0.067} & 0.288\,$\rightarrow$\,\textbf{0.135} & 0.499\,$\rightarrow$\,\textbf{0.377} & 0.687\,$\rightarrow$\,\textbf{0.359} & -- & -- \\
OpenCV-14 & 0.107\,$\rightarrow$\,\textbf{0.060} & 0.308\,$\rightarrow$\,\textbf{0.131} & 0.476\,$\rightarrow$\,\textbf{0.359} & 0.600\,$\rightarrow$\,\textbf{0.384} & 1.98\,$\rightarrow$\,\textbf{1.70} & 0.123\,$\rightarrow$\,\textbf{0.079} \\
Kannala--Brandt & 0.117\,$\rightarrow$\,\textbf{0.106} & 0.183\,$\rightarrow$\,0.176 & 0.456\,$\rightarrow$\,\textbf{0.351} & 0.567\,$\rightarrow$\,\textbf{0.329} & 2.00\,$\rightarrow$\,\textbf{1.52} & 0.113\,$\rightarrow$\,\textbf{0.063} \\
Double Sphere & 0.261\,$\rightarrow$\,\textit{0.286} & 0.219\,$\rightarrow$\,\textit{0.268} & 1.42\,$\rightarrow$\,1.34 & 0.734\,$\rightarrow$\,\textbf{0.461} & 2.09\,$\rightarrow$\,\textbf{1.49} & 0.109\,$\rightarrow$\,\textbf{0.067} \\
Central spline & 0.040\,$\rightarrow$\,\textit{0.048} & 0.161\,$\rightarrow$\,\textbf{0.136} & 0.479\,$\rightarrow$\,\textbf{0.331} & 0.504\,$\rightarrow$\,\textbf{0.366} & 1.99\,$\rightarrow$\,\textbf{1.45} & 0.121\,$\rightarrow$\,\textbf{0.100} \\
\midrule
\multicolumn{7}{@{}l}{\underline{\textit{95th percentile reprojection error [px]}}} \\
Brown--Conrady & 0.356\,$\rightarrow$\,\textbf{0.173} & 0.855\,$\rightarrow$\,\textbf{0.321} & 2.07\,$\rightarrow$\,\textbf{1.78} & 2.73\,$\rightarrow$\,\textbf{1.72} & -- & -- \\
OpenCV-14 & 0.373\,$\rightarrow$\,\textbf{0.155} & 0.969\,$\rightarrow$\,\textbf{0.337} & 2.16\,$\rightarrow$\,\textbf{1.74} & 2.87\,$\rightarrow$\,\textbf{1.93} & 9.15\,$\rightarrow$\,\textbf{5.68} & 0.818\,$\rightarrow$\,\textbf{0.440} \\
Kannala--Brandt & 0.446\,$\rightarrow$\,\textbf{0.285} & 0.730\,$\rightarrow$\,\textbf{0.515} & 1.99\,$\rightarrow$\,\textbf{1.79} & 2.37\,$\rightarrow$\,\textbf{1.75} & 6.80\,$\rightarrow$\,\textbf{5.43} & 0.590\,$\rightarrow$\,\textbf{0.254} \\
Double Sphere & 1.70\,$\rightarrow$\,\textbf{1.16} & 0.717\,$\rightarrow$\,\textbf{0.625} & 5.13\,$\rightarrow$\,\textbf{4.42} & 2.42\,$\rightarrow$\,\textbf{1.94} & 7.41\,$\rightarrow$\,\textbf{5.45} & 0.591\,$\rightarrow$\,\textbf{0.289} \\
Central spline & 0.131\,$\rightarrow$\,\textbf{0.119} & 0.436\,$\rightarrow$\,\textbf{0.331} & 2.45\,$\rightarrow$\,\textbf{1.76} & 2.35\,$\rightarrow$\,\textbf{1.97} & 7.58\,$\rightarrow$\,\textbf{5.51} & 0.996\,$\rightarrow$\,\textbf{0.538} \\
\bottomrule
\end{tabular*}

        \caption{Reprojection error on the held-out high-precision target, geometric start $\rightarrow$ radiometric refinement, for each camera model, lens and calibration target. The refined value is bold where it is significantly lower than the start and italic where it is significantly higher (95\,\% bootstrap interval of the difference, \cref{tab:supp-refinement-a}); the mean and the RMS are in \cref{tab:supp-refinement-b}. --: Brown--Conrady cannot represent the $180^\circ$ field of view of the fisheye.}
        \label{tab:refinement-chrome}
\end{table*}
\subsection{Detector Bias}
\label{sec:detector-bias}
Feature detectors introduce a measurement bias that degrades the
regressed camera model. We rendered 10,000 images each of the
AprilGrid \cite{furgale2013unified,olson2011apriltag} with an FOV
camera model \cite{devernay2001straight}, which has a mild
distortion, and of the Siemens-star target of Schöps \etal
\cite{10000params} with a Double Sphere model \cite{double-sphere},
which has a strong fisheye distortion, in evenly lit calibration
scenes. The second column of \cref{fig:detector-bias} shows
the displacement between the true projection of each feature and its
detection, for the AprilGrid detector of OpenCV \cite{opencv_library}
(top) and the Siemens-star detector of Schöps \etal (bottom). Both
detectors exhibit a bias that grows toward the image border. The third
column projects the features with the camera and the poses of a
calibration. In the geometric calibration, OpenCV-14 for the AprilGrid
and the central spline for the Siemens stars, the detection bias is
absorbed by both the camera and the poses, whereas the radiometric
refinement largely removes it. The fourth column substitutes the true
poses for the estimated ones, which isolates the camera model. This
column is the relevant one: the poses are discarded after the
calibration, and a downstream task depends on the camera alone.
The geometrically calibrated camera has absorbed a strong bias from
the feature detections, also for the highly flexible spline of Schöps
\etal, and our radiometric refinement removes most of it.




\subsection{Geometric Accuracy}
\label{sec:geometric-accuracy}
We lift four commonly used central parametric camera models,
Brown--Conrady \cite{Brown1966DecenteringDO}, the 14-parameter model
of OpenCV \cite{opencv_library}, Kannala--Brandt \cite{kannala-brandt}
and Double Sphere \cite{double-sphere}, and the central generic spline
model of Schöps \etal \cite{10000params} to the radiometric camera. We
calibrate them on three lenses chosen for their diverse fields of view
and distortions (\cref{tab:dataset-lenses}): a Tamron 20\,mm
wide-angle lens with barrel distortion ($94^\circ$ diagonal field of
view), a Sigma 14\,mm rectilinear wide-angle lens ($114^\circ$) and a
Sigma 8\,mm circular fisheye ($180^\circ$). We captured two
standard calibration targets, the Siemens-star target of Schöps \etal
\cite{10000params} and an AprilGrid \cite{furgale2013unified}, each
$1\,\mathrm{m}\times1\,\mathrm{m}$ and printed directly on Alu-Dibond,
at $f/22$ (the fisheye's Siemens
stars at $f/32$) and ISO\,100, to minimize defocus and thus give
optimal conditions to the geometric camera models, which do not model
defocus. Each model is started from the geometric calibration of its
native tool where the tool supports the model and the target, and
from our own estimator otherwise (\cref{sec:supp-starts}). The
radiometric refinement samples 8,192 photosites per image for the
parametric models and 16,384 for the spline.

The test set consists of backlit images of the Thorlabs R3L3S7N1
high-precision distortion target \cite{thorlabsR3L3S7N1}, captured
by rotating the camera about a pivot, which reduces the poses to one
rotation per view. Every dot is evaluated with a pose that it did
not help to fit (\cref{sec:supp-chrome-evaluation}). 
% \begin{figure}[t]
%         \centering
%         \input{fig/color-wheels}
%         \caption{\protect\input{fig/color-wheels-caption}}
%         \label{fig:color-wheels}
% \end{figure}

\Cref{tab:refinement-chrome} compares the reprojection error of each
model at its geometric start and after the radiometric refinement on
the Thorlabs test set (\cref{tab:supp-dataset}). We call a change significant where the 95\,\% bootstrap interval of
the paired difference excludes zero (\cref{sec:supp-chrome-evaluation}).
The radiometric refinement lowers the
median reprojection error significantly in 82\,\% of the 28\
combinations of model, lens and target (23, by 9 to 57\,\%), raises
it significantly in 11\,\% (three, all on the Tamron 20\,mm), and
lowers the $95$th percentile significantly in all of them, by 9 to
65\,\%. On the Sigma 14\,mm and on the fisheye, no model is
significantly worse in any statistic.
\Cref{fig:convergence-rms} follows the two refinements of the
central spline on the Sigma 14\,mm over their iterations. In both,
the error on the held-out target falls throughout, while the error
with respect to the detections that the geometric calibration was
fitted to rises after the first four iterations (from $0.59$ to
$0.71$\,px on the Siemens stars and from $0.56$ to $0.68$\,px on the
AprilGrid). The test error decreases while the training error
increases, which shows the bias: the refinement moves the camera
away from the biased detections that its start was fitted to. We attribute part of this bias to the freedom of the
geometric calibration: The feature wise parametrization of the target in \cite{10000params} allows to absorb the bias, while our lower parameter parametrization (\cref{sec:target-shape}) cannot. 

\begin{figure}[t]
        \centering
        % convergence-rms.tex — generated by convergence/tools/rms_figure.py; do not edit.
% The RMS per LM iteration of the error against the detections and of the held-out reprojection error on the chrome target, one panel
% per case (one PGFPlots picture each, its key under it), composed by a LaTeX tabular, the panels beside each other: where the line
% is at least 4.5in wide (a figure* at text width) the panels convergence-rms/<case>.tex, in one column (a figure) their narrow
% renderings convergence-rms/<case>-narrow.tex (title on two lines). The key: the red curve's entry under each panel (the panel's own), the
% blue curve's entry once, in the figure's centre (at text width between the red entries, in one column on a line under them).
% Needs pgfplots (compat 1.18). \rcfigdir is the folder this file lies in, as the host document sees it (default fig).
\providecommand{\rcfigdir}{fig}%
\begingroup
\ifdefined\rcpanelwd\else\newlength\rcpanelwd\fi\def\rcpanelwidth{\rcpanelwd}%
\setlength{\tabcolsep}{0pt}\definecolor{rcchr}{HTML}{0000FF}\definecolor{rcapr}{HTML}{FF0000}\definecolor{rcsie}{HTML}{228833}%
\def\rckey#1#2{\textcolor{#1}{\rule[0.5ex]{0.16in}{0.9pt}}\hspace{0.04in}#2}%
\ifdim\linewidth<4.5in\relax
  \setlength{\rcpanelwd}{\dimexpr(\linewidth-62pt)/2\relax}%
  \begin{tabular}{@{}l@{\hspace{4pt}}l@{}}
  {\def\rcpanelylabel{Median error [px]}\hspace*{31pt}\input{\rcfigdir/convergence-rms/si30_spline_warped-narrow}\hspace*{4pt}} &
  {\def\rcpanelylabel{}\hspace*{19pt}\input{\rcfigdir/convergence-rms/ag30_spline_warped-narrow}\hspace*{4pt}} \\[-9pt]
  \multicolumn{2}{@{}c@{}}{\scriptsize \rckey{rcapr}{AprilGrid detections}\hspace{0.1in}\rckey{rcsie}{Siemens-star detections}\hspace{0.1in}\rckey{rcchr}{high-precision target}} \\
  \end{tabular}%
\else
  \setlength{\rcpanelwd}{\dimexpr(\linewidth-80pt)/2\relax}%
  \begin{tabular}{@{}l@{\hspace{10pt}}l@{}}
  {\def\rcpanelylabel{Median reprojection error [px]}\hspace*{36pt}\input{\rcfigdir/convergence-rms/si30_spline_warped}\hspace*{6pt}} &
  {\def\rcpanelylabel{}\hspace*{22pt}\input{\rcfigdir/convergence-rms/ag30_spline_warped}\hspace*{6pt}} \\[3pt]
  \multicolumn{2}{@{}c@{}}{\scriptsize \rckey{rcapr}{AprilGrid detections}\hspace{0.15in}\rckey{rcsie}{Siemens-star detections}\hspace{0.15in}\rckey{rcchr}{high-precision target}} \\
  \end{tabular}%
\fi
\endgroup

        \caption{Median reprojection error during the radiometric refinement of the central spline on the Sigma 14\,mm, with respect to the detections that the geometric calibration was fitted to and on the held-out high-precision target. After the first iterations the curves diverge in both panels: the error on the detections rises while the error on the held-out target falls, which shows that the geometric calibration had absorbed the bias of its detector.}
        \label{fig:convergence-rms}
\end{figure}

\subsection{Vignetting}
\label{sec:experiment-vignetting}
\begin{figure}[t]
        \centering
        % vignetting.tex — generated by experiments/pgf/vignetting.py; do not edit.
% The figure: its 3 panels (one PGFPlots picture each, vignetting/*.tex) stacked in a tabular; the panels of a lens share the x axis.
% The column titles, the row names, the axis name and the legend are LaTeX text in the host document's fonts. \input it inside a
% figure; it needs \usepackage{pgfplots} (which loads graphicx), \pgfplotsset{compat=1.18} and amssymb. \rcfigdir is the folder this file lies in, as the
% host document sees it (default fig).
\providecommand{\rcfigdir}{fig}%
\begingroup
\ifdefined\rcpanelwd\else\newlength\rcpanelwd\fi
\setlength{\rcpanelwd}{\dimexpr(\linewidth-52pt)/3\relax}\def\rcpanelwidth{\rcpanelwd}%
\definecolor{rcvlearned}{HTML}{0000FF}\definecolor{rcvflat}{HTML}{00A000}\definecolor{rcvthin}{HTML}{FF0000}%
\footnotesize\setlength{\tabcolsep}{0pt}%
\begin{tabular}{@{}c@{\hspace{24pt}}c@{\hspace{7pt}}c@{\hspace{7pt}}c@{\hspace{4pt}}}
  & \makebox[0pt]{\scriptsize\begin{tabular}[b]{@{}c@{}}Sigma 14\,mm lens\\$f/4$\end{tabular}} & \makebox[0pt]{\scriptsize\begin{tabular}[b]{@{}c@{}}Fisheye lens\\$f/4$\end{tabular}} & \makebox[0pt]{\scriptsize\begin{tabular}[b]{@{}c@{}}Tamron 20\,mm lens\\$f/2.8$\end{tabular}} \\[2pt]
  \smash{\raisebox{-0.5\height}{\rotatebox{90}{$V / V(\mathfrak{c})$}}} &
  \input{\rcfigdir/vignetting/14mm-diagonal} &
  \pgfplotsset{every axis/.append style={yticklabels=\empty}}\input{\rcfigdir/vignetting/fisheye-diagonal} &
  \pgfplotsset{every axis/.append style={yticklabels=\empty}}\input{\rcfigdir/vignetting/tamron-diagonal} \\\noalign{\vskip 1pt}
  & \multicolumn{3}{@{}c@{\hspace{4pt}}}{position along the diagonal [px]} \\\noalign{\vskip 5pt}
  % the legend: one line, shrunk only when it is wider than the line
  \multicolumn{4}{@{}c@{}}{\resizebox{\ifdim\width>\linewidth\linewidth\else\width\fi}{!}{\mbox{\tikz[baseline=-0.6ex]{\draw[rcvflat, line width=1pt] (0,0) -- (12pt,0);}\,flat field}\enspace \mbox{\tikz[baseline=-0.6ex]{\draw[rcvlearned, line width=1pt, dash pattern=on 3pt off 1.5pt] (0,0) -- (12pt,0);}\,learned (ours)}\enspace \mbox{\tikz[baseline=-0.6ex]{\draw[rcvthin, line width=1pt] (0,0) -- (12pt,0);}\,thin lens model ($\cos^4\alpha$)}}} \\
\end{tabular}%
\endgroup

        \caption{Normalized vignetting along the direction of the sensor diagonal through the principal point: the $\cos^4$ falloff of the thin lens and our learned vignetting against the measured flat field.}
        \label{fig:vignetting}
\end{figure}
% Figure of Sec. 5.4 (Defocus), declared here so that it floats to the top of that section's page.
\begin{figure*}[t]
        \centering
        \begin{minipage}[c]{0.43\textwidth}
                \centering
                % focus-surfaces.tex — generated by experiments/pgf/focus_surfaces.py; do not edit.
% The figure: its 2 panels (one PGFPlots picture each, a 3-D axis, focus-surfaces/*.tex) side by side in a tabular, their axis boxes at one
% height, under each its key (the cells' colour: focus distance, light = near, dark = far) and under the keys their name. The titles, the keys' numbers and their name are LaTeX text in the host document's fonts, and so are the axes' names and numbers inside the panels. \input it inside a
% figure (one column); it needs \usepackage{pgfplots} and \pgfplotsset{compat=1.18}. \rcfigdir is the folder this file lies in, as the host
% document sees it (default fig).
\providecommand{\rcfigdir}{fig}%
\begingroup
\footnotesize\setlength{\tabcolsep}{0pt}%
\begin{tabular}{@{\hspace{22pt}}c@{\hspace{34pt}}c@{\hspace{20pt}}}
  \makebox[0pt]{Fisheye lens, $f/4$} & \makebox[0pt]{Tamron 20\,mm lens, $f/8$} \\\noalign{\vskip 3pt}
  \def\rcpanelwidth{\dimexpr(\linewidth-76pt)*563/1000\relax}\input{\rcfigdir/focus-surfaces/fisheye} &
  \def\rcpanelwidth{\dimexpr(\linewidth-76pt)*437/1000\relax}\input{\rcfigdir/focus-surfaces/tamron} \\\noalign{\vskip 3pt}
  \makebox[0pt]{{\scriptsize 0.52\,m\,\tikz[baseline=0.5pt]{\fill[fill={rgb,255:red,230;green,239;blue,249}] (0pt,0) rectangle (2.05pt,4pt);\fill[fill={rgb,255:red,222;green,234;blue,246}] (2pt,0) rectangle (4.05pt,4pt);\fill[fill={rgb,255:red,214;green,228;blue,243}] (4pt,0) rectangle (6.05pt,4pt);\fill[fill={rgb,255:red,206;green,222;blue,240}] (6pt,0) rectangle (8.05pt,4pt);\fill[fill={rgb,255:red,198;green,217;blue,237}] (8pt,0) rectangle (10.05pt,4pt);\fill[fill={rgb,255:red,190;green,211;blue,234}] (10pt,0) rectangle (12.05pt,4pt);\fill[fill={rgb,255:red,182;green,206;blue,231}] (12pt,0) rectangle (14.05pt,4pt);\fill[fill={rgb,255:red,174;green,200;blue,228}] (14pt,0) rectangle (16.05pt,4pt);\fill[fill={rgb,255:red,166;green,195;blue,224}] (16pt,0) rectangle (18.05pt,4pt);\fill[fill={rgb,255:red,158;green,189;blue,221}] (18pt,0) rectangle (20.05pt,4pt);\fill[fill={rgb,255:red,150;green,183;blue,218}] (20pt,0) rectangle (22.05pt,4pt);\fill[fill={rgb,255:red,142;green,178;blue,215}] (22pt,0) rectangle (24.05pt,4pt);\fill[fill={rgb,255:red,135;green,172;blue,212}] (24pt,0) rectangle (26.05pt,4pt);\fill[fill={rgb,255:red,127;green,167;blue,209}] (26pt,0) rectangle (28.05pt,4pt);\fill[fill={rgb,255:red,119;green,161;blue,206}] (28pt,0) rectangle (30.05pt,4pt);\fill[fill={rgb,255:red,111;green,155;blue,203}] (30pt,0) rectangle (32.05pt,4pt);\fill[fill={rgb,255:red,103;green,150;blue,199}] (32pt,0) rectangle (34.05pt,4pt);\fill[fill={rgb,255:red,95;green,144;blue,196}] (34pt,0) rectangle (36.05pt,4pt);\fill[fill={rgb,255:red,87;green,139;blue,193}] (36pt,0) rectangle (38.05pt,4pt);\fill[fill={rgb,255:red,79;green,133;blue,190}] (38pt,0) rectangle (40.05pt,4pt);\fill[fill={rgb,255:red,71;green,128;blue,187}] (40pt,0) rectangle (42.05pt,4pt);\fill[fill={rgb,255:red,63;green,122;blue,184}] (42pt,0) rectangle (44.05pt,4pt);\fill[fill={rgb,255:red,55;green,116;blue,181}] (44pt,0) rectangle (46.05pt,4pt);\fill[fill={rgb,255:red,47;green,111;blue,178}] (46pt,0) rectangle (48.05pt,4pt);\draw[line width=0.2pt] (0,0) rectangle (48pt,4pt);}\,0.76\,m}} & \makebox[0pt]{{\scriptsize 0.22\,m\,\tikz[baseline=0.5pt]{\fill[fill={rgb,255:red,230;green,239;blue,249}] (0pt,0) rectangle (2.05pt,4pt);\fill[fill={rgb,255:red,222;green,234;blue,246}] (2pt,0) rectangle (4.05pt,4pt);\fill[fill={rgb,255:red,214;green,228;blue,243}] (4pt,0) rectangle (6.05pt,4pt);\fill[fill={rgb,255:red,206;green,222;blue,240}] (6pt,0) rectangle (8.05pt,4pt);\fill[fill={rgb,255:red,198;green,217;blue,237}] (8pt,0) rectangle (10.05pt,4pt);\fill[fill={rgb,255:red,190;green,211;blue,234}] (10pt,0) rectangle (12.05pt,4pt);\fill[fill={rgb,255:red,182;green,206;blue,231}] (12pt,0) rectangle (14.05pt,4pt);\fill[fill={rgb,255:red,174;green,200;blue,228}] (14pt,0) rectangle (16.05pt,4pt);\fill[fill={rgb,255:red,166;green,195;blue,224}] (16pt,0) rectangle (18.05pt,4pt);\fill[fill={rgb,255:red,158;green,189;blue,221}] (18pt,0) rectangle (20.05pt,4pt);\fill[fill={rgb,255:red,150;green,183;blue,218}] (20pt,0) rectangle (22.05pt,4pt);\fill[fill={rgb,255:red,142;green,178;blue,215}] (22pt,0) rectangle (24.05pt,4pt);\fill[fill={rgb,255:red,135;green,172;blue,212}] (24pt,0) rectangle (26.05pt,4pt);\fill[fill={rgb,255:red,127;green,167;blue,209}] (26pt,0) rectangle (28.05pt,4pt);\fill[fill={rgb,255:red,119;green,161;blue,206}] (28pt,0) rectangle (30.05pt,4pt);\fill[fill={rgb,255:red,111;green,155;blue,203}] (30pt,0) rectangle (32.05pt,4pt);\fill[fill={rgb,255:red,103;green,150;blue,199}] (32pt,0) rectangle (34.05pt,4pt);\fill[fill={rgb,255:red,95;green,144;blue,196}] (34pt,0) rectangle (36.05pt,4pt);\fill[fill={rgb,255:red,87;green,139;blue,193}] (36pt,0) rectangle (38.05pt,4pt);\fill[fill={rgb,255:red,79;green,133;blue,190}] (38pt,0) rectangle (40.05pt,4pt);\fill[fill={rgb,255:red,71;green,128;blue,187}] (40pt,0) rectangle (42.05pt,4pt);\fill[fill={rgb,255:red,63;green,122;blue,184}] (42pt,0) rectangle (44.05pt,4pt);\fill[fill={rgb,255:red,55;green,116;blue,181}] (44pt,0) rectangle (46.05pt,4pt);\fill[fill={rgb,255:red,47;green,111;blue,178}] (46pt,0) rectangle (48.05pt,4pt);\draw[line width=0.2pt] (0,0) rectangle (48pt,4pt);}\,0.31\,m}} \\\noalign{\vskip 1pt}
  \multicolumn{2}{@{\hspace{22pt}}c@{\hspace{20pt}}}{\scriptsize focus distance from the camera centre} \\
\end{tabular}%
\endgroup

        \end{minipage}\hfill
        \begin{minipage}[c]{0.54\textwidth}
                \centering
                % texture-learned.tex — generated by texture-learned/tools/texture_learned.py; do not edit.
% The figure: its 4 bands (one PGFPlots picture each, texture-learned/*.tex) one under the other in a tabular, without a gap: the same text lines of the page,
% once from each source. The bands' names are LaTeX text in the host document's
% fonts, in a margin at the left; the short rules between them separate the names. \input it inside a figure (one column); it needs
% \usepackage{pgfplots} and \pgfplotsset{compat=1.18}. \rcfigdir is the folder this file lies in, as the host document sees it (default
% fig). The bands are as wide as the line less the names' margin (44 pt) and the gap (4 pt).
\providecommand{\rcfigdir}{fig}%
\begingroup
\ifdefined\rcpanelwd\else\newlength\rcpanelwd\fi
\setlength{\rcpanelwd}{\dimexpr\linewidth-48pt\relax}\def\rcpanelwidth{\rcpanelwd}%
\footnotesize\setlength{\tabcolsep}{0pt}%
\begin{tabular}{@{}c@{\hspace{4pt}}c@{}}
  \tikz[baseline=(n.center)]{\node[inner sep=0pt, outer sep=0pt, minimum width=44pt, minimum height=0.13078\rcpanelwd, align=center, font=\fontsize{6.5}{7.5}\selectfont] (n) {ground\\truth};\draw[overlay, black!60, line width=0.4pt] (n.south west) -- (n.south east);} &
  \input{\rcfigdir/texture-learned/truth} \\
  \tikz[baseline=(n.center)]{\node[inner sep=0pt, outer sep=0pt, minimum width=44pt, minimum height=0.13264\rcpanelwd, align=center, font=\fontsize{6.5}{7.5}\selectfont] (n) {training\\image\\39 of 63};\draw[overlay, black!60, line width=0.4pt] (n.south west) -- (n.south east);} &
  \input{\rcfigdir/texture-learned/photograph} \\
  \tikz[baseline=(n.center)]{\node[inner sep=0pt, outer sep=0pt, minimum width=44pt, minimum height=0.13078\rcpanelwd, align=center, font=\fontsize{6.5}{7.5}\selectfont] (n) {pinhole};\draw[overlay, black!60, line width=0.4pt] (n.south west) -- (n.south east);} &
  \input{\rcfigdir/texture-learned/pinhole} \\
  \tikz[baseline=(n.center)]{\node[inner sep=0pt, outer sep=0pt, minimum width=44pt, minimum height=0.13078\rcpanelwd, align=center, font=\fontsize{6.5}{7.5}\selectfont] (n) {ours};} &
  \input{\rcfigdir/texture-learned/ours} \\
\end{tabular}%
\endgroup

        \end{minipage}
        \caption{The focus model. \emph{Left:} the learned focus surfaces, a hemisphere for the fisheye at $f/4$ and a plane for the Tamron 20\,mm at $f/8$. \emph{Right:} a texture estimated from defocused images. The radiometric camera recovers sharp text, whereas the pinhole camera has to put the blur of the images into the texture.}
        \label{fig:focus-surfaces}
        \label{fig:texture-learned}
\end{figure*}
We use the three lenses of \cref{sec:geometric-accuracy} and the
spline model with Siemens-star captures at a large aperture ($f/4$; $f/2.8$ for the
Tamron) for stronger vignetting, and one flat-field image per lens
(\cref{tab:supp-dataset}): the lens, with its hood on, covered by a
white acrylic plate and lit by diffuse daylight through a light
diffuser. \Cref{fig:vignetting} compares, along the
direction of the sensor diagonal through the principal point, the
learned vignetting, a B-spline (\cref{sec:vignetting}),
with the flat field, read off the green photosites, and with the
$\cos^4$ falloff of the thin lens. All profiles are normalized at the
principal point. The learned vignetting follows the flat field for all
three lenses: its RMS deviation along this line, in units of the
value at the principal point, is $0.009$ for the 14\,mm lens, $0.042$ for the
fisheye and $0.025$ for the Tamron, against $0.244$, $0.199$ and
$0.059$ for the $\cos^4$ falloff.


\subsection{Defocus}
\label{sec:experiment-defocus}
We validate the focus model in three ways. First,
\cref{fig:focus-surfaces} (left) shows the focus surfaces learned with the
spline model for the Tamron 20\,mm at $f/8$ and for the fisheye at
$f/4$: a plane for the wide-angle lens and a hemisphere for the
fisheye, as expected for these lens types
\cite{Kingslake2010,wide-field-of-view-2017}.

Second, we solve a first multi-view radiometric problem with the
calibrated camera: we estimate the unknown reflectance of a surface
from many defocused views of it by inverse rendering. With a pinhole
aperture, the estimated texture has to explain the defocus of the
images and is blurred. We calibrate the Sigma 14\,mm at $f/2.8$ on
Siemens stars, once with the aperture and the focus surface and once
as a pinhole, and estimate the texture of a board with printed text
from 63 defocused images with each camera. \Cref{fig:texture-learned} (right)
compares both textures with the printed design: the radiometric
camera recovers sharp text where the pinhole camera returns the blur
of the images. Since the calibration works with relative irradiance,
a texture is recovered only up to a scale and an offset, which we fit
once per texture before comparing it with the design. With the design
ranging from 0 to 1, the RMS difference to it is $0.177$ for the
radiometric camera and $0.235$ for the pinhole, and it is lower for
the radiometric camera on 99\,\% of the page.




Third, we ablate the focus model on two lenses with visible defocus,
the Tamron 20\,mm at $f/8$ with the spline model and a 24--70\,mm zoom
lens at 48\,mm and $f/10$ with Brown--Conrady, by refining the same
geometric start once with the aperture and the focus
surface and once as a pinhole. \Cref{tab:defocus} shows that modeling
the aperture and the focus surface lowers the median reprojection
error on the Thorlabs target by 25\,\% and 16\,\% compared with the
pinhole refinement, both significantly; on the Tamron, the pinhole
refinement is even worse than the geometric start.
\begin{table}[t]
        \centering
        \footnotesize
\setlength{\tabcolsep}{1.5pt}
\begin{tabular*}{\linewidth}{@{\extracolsep{\fill}}lll@{}}
\toprule
\textbf{Refinement} & \textbf{Median [px]} & \textbf{95th perc. [px]} \\
\midrule
\multicolumn{3}{@{}l}{\underline{\textit{Tamron 20\,mm, $f/8$, central spline}}} \\
Geometric start & 0.160 & 0.499 \\
Pinhole & \textit{0.199} & \textit{0.534} \\
Aperture and focus & \textbf{0.149} & \textbf{0.410} \\
\quad $-$ pinhole & $-0.051$ {\scriptsize[$-0.054$,\,$-0.047$]} & $-0.124$ {\scriptsize[$-0.139$,\,$-0.103$]} \\
\midrule
\multicolumn{3}{@{}l}{\underline{\textit{Zoom at 48\,mm, $f/10$, Brown--Conrady}}} \\
Geometric start & 0.068 & 0.194 \\
Pinhole & \textbf{0.050} & \textbf{0.150} \\
Aperture and focus & \textbf{0.042} & \textbf{0.117} \\
\quad $-$ pinhole & $-0.008$ {\scriptsize[$-0.013$,\,$-0.003$]} & $-0.032$ {\scriptsize[$-0.047$,\,$-0.017$]} \\
\bottomrule
\end{tabular*}

        \caption{Defocus ablation on the held-out target: the refinement as a pinhole and with the aperture and the focus surface, and their difference with its 95\,\% bootstrap interval. Bold and italic as in \cref{tab:refinement-chrome}.}
        \label{tab:defocus}
\end{table}
